% GENERATED FILE. Edit canonical lessons, then run npm run generate:books.
% canonical-source-hash: fnv1a64:54463a7730f9c84c
% canonical-lessons: KA-C488-bisigali, KA-C488-uribisilu, KA-C488-battu, KA-C488-biruku, KA-C488-taggu

\chapter{Heat and Dry Spells}
\label{ch:ka-heat-and-dry-spells}
\hlchaptermodality{488}

\begin{chapteropening}
\textbf{By the end of this chapter:} \emph{I can say a heatwave, scorching sun, to run dry, a crack and to subside.}

Complete the last lesson of chapter 488: I can say a heatwave, scorching sun, to run dry, a crack and to subside.
\end{chapteropening}

\section[\texorpdfstring{\kn{ಬಿಸಿಗಾಳಿ} (bisigāḷi)}{ಬಿಸಿಗಾಳಿ (bisigāḷi)}]{\kn{ಬಿಸಿಗಾಳಿ} (bisigāḷi) — a heatwave; a hot wind}

\label{lesson:KA-C488-bisigali}



\begin{quote}
Before the new one: say the Kannada for a landslide, then the Kannada for a tsunami.
\end{quote}

\subsection*{You'll want to know: \kn{ಬಿಸಿಗಾಳಿ}}
\textbf{\kn{ಬಿಸಿಗಾಳಿ}} — \emph{bisigāḷi} — \textquotedblleft{}a heatwave; a hot wind\textquotedblright{}.

\begin{quote}
\kn{ಬಿಸಿಗಾಳಿಯ} \kn{ಕಾರಣ} \kn{ಮಧ್ಯಾಹ್ನ} \kn{ಹೊರಗೆ} \kn{ಹೋಗಬಾರದು}.
\end{quote}

\begin{quote}
\emph{bisigāḷiya kāraṇa madhyāhna horage hōgabāradu.}
\end{quote}

\begin{quote}
Because of the heatwave nobody should go out at midday.
\end{quote}

\kn{ಬಿಸಿ}, hot, plus \kn{ಗಾಳಿ}, wind.

\begin{cousinweb}
\textbf{\kn{ಬಿಸಿಗಾಳಿ}} (\emph{bisigāḷi}) is a native Kannada compound of \emph{bisi}, \textquotedblleft{}heat, hot\textquotedblright{} (DEDR 5517), and \kn{ಗಾಳಿ}, \textquotedblleft{}wind\textquotedblright{} (DEDR 1499).
\end{cousinweb}

\begin{grammarlens}[title={What you've built}]
One more word for getting around.
\end{grammarlens}

\subsection*{Your turn}
\begin{itemize}
\raggedright
  \item \emph{Say it:} \emph{bisigāḷi}
  \item \emph{Say it:} \emph{bisigāḷi}, once more
  \item \emph{Say it:} \emph{bisigāḷiya kāraṇa madhyāhna horage hōgabāradu.} — Because of the heatwave nobody should go out at midday.
  \item \emph{Recall:} say the Kannada for a landslide, then the Kannada for a tsunami, then say \emph{bisigāḷi} again
\end{itemize}

\begin{tcolorbox}[breakable,colback=teal!4,colframe=teal!35!black,title={Before you move on}]
What does \kn{ಬಿಸಿಗಾಳಿ} mean? (\textquotedblleft{}A heatwave; a hot wind\textquotedblright{}.) Say it once more.
\end{tcolorbox}
\section[\texorpdfstring{\kn{ಉರಿಬಿಸಿಲು} (uribisilu)}{ಉರಿಬಿಸಿಲು (uribisilu)}]{\kn{ಉರಿಬಿಸಿಲು} (uribisilu) — scorching sun}

\label{lesson:KA-C488-uribisilu}



\begin{quote}
Before the new one: say the Kannada for a tsunami, then the Kannada for a heatwave.
\end{quote}

\subsection*{You'll want to know: \kn{ಉರಿಬಿಸಿಲು}}
\textbf{\kn{ಉರಿಬಿಸಿಲು}} — \emph{uribisilu} — \textquotedblleft{}scorching sun\textquotedblright{}.

\begin{quote}
\kn{ಉರಿಬಿಸಿಲಿನಲ್ಲಿ} \kn{ಕೂಲಿಗಳು} \kn{ರಸ್ತೆ} \kn{ಕೆಲಸ} \kn{ಮಾಡುತ್ತಿದ್ದರು}.
\end{quote}

\begin{quote}
\emph{uribisilinalli kūligaḷu raste kelasa māḍuttiddaru.}
\end{quote}

\begin{quote}
Labourers were working on the road in the scorching sun.
\end{quote}

Burning sunshine: the hottest hours of a summer day.

\begin{cousinweb}
\textbf{\kn{ಉರಿಬಿಸಿಲು}} (\emph{uribisilu}) is a native Kannada compound of \emph{uri}, \textquotedblleft{}to burn, to blaze\textquotedblright{} (DEDR 656), and \kn{ಬಿಸಿಲು}, \textquotedblleft{}the heat and glare of the sun\textquotedblright{} (DEDR 5517).
\end{cousinweb}

\begin{grammarlens}[title={What you've built}]
One more word for getting around.
\end{grammarlens}

\subsection*{Your turn}
\begin{itemize}
\raggedright
  \item \emph{Say it:} \emph{uribisilu}
  \item \emph{Say it:} \emph{uribisilu}, once more
  \item \emph{Say it:} \emph{uribisilinalli kūligaḷu raste kelasa māḍuttiddaru.} — Labourers were working on the road in the scorching sun.
  \item \emph{Recall:} say the Kannada for a tsunami, then the Kannada for a heatwave, then say \emph{uribisilu} again
\end{itemize}

\begin{tcolorbox}[breakable,colback=teal!4,colframe=teal!35!black,title={Before you move on}]
What does \kn{ಉರಿಬಿಸಿಲು} mean? (\textquotedblleft{}Scorching sun\textquotedblright{}.) Say it once more.
\end{tcolorbox}
\section[\texorpdfstring{\kn{ಬತ್ತು} (battu)}{ಬತ್ತು (battu)}]{\kn{ಬತ್ತು} (battu) — to run dry; to dry up}

\label{lesson:KA-C488-battu}



\begin{quote}
Before the new one: say the Kannada for a heatwave, then the Kannada for scorching sun.
\end{quote}

\subsection*{You'll want to know: \kn{ಬತ್ತು}}
\textbf{\kn{ಬತ್ತು}} — \emph{battu} — \textquotedblleft{}to run dry; to dry up\textquotedblright{}.

\begin{quote}
\kn{ಬರದಿಂದ} \kn{ಊರಿನ} \kn{ಕೆರೆ} \kn{ಬತ್ತಿತು}.
\end{quote}

\begin{quote}
\emph{baradinda ūrina kere battitu.}
\end{quote}

\begin{quote}
Because of the drought the town lake ran dry.
\end{quote}

Past (it …): \kn{ಬತ್ತಿತು}, \emph{battitu}.

Of wells, lakes, rivers and springs.

\begin{cousinweb}
\textbf{\kn{ಬತ್ತು}} (\emph{battu}) is native Dravidian: DEDR 5320 lists Kannada \emph{battu}, \textquotedblleft{}to grow dry as a stream, to dry up, to evaporate\textquotedblright{}, in the family of \kn{ಬರ}, a drought.
\end{cousinweb}

\begin{grammarlens}[title={What you've built}]
One more word for getting around.
\end{grammarlens}

\subsection*{Your turn}
\begin{itemize}
\raggedright
  \item \emph{Say it:} \emph{battu}
  \item \emph{Say it:} \emph{battu}, once more
  \item \emph{Say it:} \emph{baradinda ūrina kere battitu.} — Because of the drought the town lake ran dry.
  \item \emph{Recall:} say the Kannada for a heatwave, then the Kannada for scorching sun, then say \emph{battu} again
\end{itemize}

\begin{tcolorbox}[breakable,colback=teal!4,colframe=teal!35!black,title={Before you move on}]
What does \kn{ಬತ್ತು} mean? (\textquotedblleft{}To run dry; to dry up\textquotedblright{}.) Say it once more.
\end{tcolorbox}
\section[\texorpdfstring{\kn{ಬಿರುಕು} (biruku)}{ಬಿರುಕು (biruku)}]{\kn{ಬಿರುಕು} (biruku) — a crack; a fissure}

\label{lesson:KA-C488-biruku}



\begin{quote}
Before the new one: say the Kannada for scorching sun, then the Kannada for to run dry.
\end{quote}

\subsection*{You'll want to know: \kn{ಬಿರುಕು}}
\textbf{\kn{ಬಿರುಕು}} — \emph{biruku} — \textquotedblleft{}a crack; a fissure\textquotedblright{}.

\begin{quote}
\kn{ಬರಗಾಲದಲ್ಲಿ} \kn{ಹೊಲದ} \kn{ಮಣ್ಣಿನಲ್ಲಿ} \kn{ಬಿರುಕುಗಳು} \kn{ಬಂದವು}.
\end{quote}

\begin{quote}
\emph{baragāladalli holada maṇṇinalli birukugaḷu bandavu.}
\end{quote}

\begin{quote}
In the drought cracks opened in the soil of the fields.
\end{quote}

Plural: \kn{ಬಿರುಕುಗಳು}, \emph{birukugaḷu}.

\kn{ಬಿರುಕು} \kn{ಬಿಡು}, to crack.

\begin{cousinweb}
\textbf{\kn{ಬಿರುಕು}} (\emph{biruku}) is native Dravidian: DEDR 5411 lists Kannada \emph{biriku}, \emph{biruku}, \textquotedblleft{}a cleft, a fissure\textquotedblright{}, beside \emph{biri}, \textquotedblleft{}to burst open\textquotedblright{}.
\end{cousinweb}

\begin{grammarlens}[title={What you've built}]
One more word for getting around.
\end{grammarlens}

\subsection*{Your turn}
\begin{itemize}
\raggedright
  \item \emph{Say it:} \emph{biruku}
  \item \emph{Say it:} \emph{biruku}, once more
  \item \emph{Say it:} \emph{baragāladalli holada maṇṇinalli birukugaḷu bandavu.} — In the drought cracks opened in the soil of the fields.
  \item \emph{Recall:} say the Kannada for scorching sun, then the Kannada for to run dry, then say \emph{biruku} again
\end{itemize}

\begin{tcolorbox}[breakable,colback=teal!4,colframe=teal!35!black,title={Before you move on}]
What does \kn{ಬಿರುಕು} mean? (\textquotedblleft{}A crack; a fissure\textquotedblright{}.) Say it once more.
\end{tcolorbox}
\section[\texorpdfstring{\kn{ತಗ್ಗು} (taggu)}{ತಗ್ಗು (taggu)}]{\kn{ತಗ್ಗು} (taggu) — to subside; to go down}

\label{lesson:KA-C488-taggu}



\begin{quote}
Before the new one: say the Kannada for to run dry, then the Kannada for a crack.
\end{quote}

\subsection*{You'll want to know: \kn{ತಗ್ಗು}}
\textbf{\kn{ತಗ್ಗು}} — \emph{taggu} — \textquotedblleft{}to subside; to go down\textquotedblright{}.

\begin{quote}
\kn{ಸಂಜೆ} \kn{ಮಳೆ} \kn{ಬಂದ} \kn{ಮೇಲೆ} \kn{ಸೆಕೆ} \kn{ತಗ್ಗಿತು}.
\end{quote}

\begin{quote}
\emph{sañje maḷe banda mēle seke taggitu.}
\end{quote}

\begin{quote}
After it rained in the evening the heat died down.
\end{quote}

Past (it …): \kn{ತಗ್ಗಿತು}, \emph{taggitu}.

As a noun, a low spot or a hollow: \kn{ರಸ್ತೆಯ} \kn{ತಗ್ಗು}.

\begin{cousinweb}
\textbf{\kn{ತಗ್ಗು}} (\emph{taggu}) is native Dravidian: DEDR 3178 lists Kannada \emph{targu}, \textquotedblleft{}to become depressed, low or less, to decline\textquotedblright{}.
\end{cousinweb}

\begin{grammarlens}[title={What you've built}]
One more word for getting around.
\end{grammarlens}

\subsection*{Your turn}
\begin{itemize}
\raggedright
  \item \emph{Say it:} \emph{taggu}
  \item \emph{Say it:} \emph{taggu}, once more
  \item \emph{Say it:} \emph{sañje maḷe banda mēle seke taggitu.} — After it rained in the evening the heat died down.
  \item \emph{Recall:} say the Kannada for to run dry, then the Kannada for a crack, then say \emph{taggu} again
\end{itemize}

\begin{tcolorbox}[breakable,colback=teal!4,colframe=teal!35!black,title={Before you move on}]
What does \kn{ತಗ್ಗು} mean? (\textquotedblleft{}To subside; to go down\textquotedblright{}.) Say it once more.
\end{tcolorbox}
